%------------------------------------------------------------------------------%
% Luis A. D'Afonseca
%
%------------------------------------------------------------------------------%
\documentclass[a4paper,12pt,fleqn]{article}

\usepackage{style_avaliacao}

\ShowAnswers

%------------------------------------------------------------------------------%
\begin{document}

\makeHeader{Integrais e Séries}{Exame -- Turma 8}{25/02/2025}

\vspace{\baselineskip}
Tabela de integrais
\begin{align*}
  \int \frac{dx}{a^2+x^2} & = \frac{1}{a}\arctan\left(\frac{x}{a}\right) + c &
  \int \csc^2(x) \, dx & = -\cot(x) + c
\end{align*}


% 01
%------------------------------------------------------------------------------%
%\vspace{\baselineskip}
\exercise{20}
Encontre o volume do sólido obtido pela rotação do triângulo de vértices
\((0, 0)\), \((1, 1)\) e \((2, -1)\) em torno da reta \(y=-1\)

\begin{answer}
  Dividimos o sólido em duas partes com $0\leq x\leq 1$ e $1<x\leq 2$

  Em ambas as partes vamos calcular o volume por cascas cilíndricas
  \[
    V = \int_a^b 2\pi r(x) h(x) \, dx
  \]

  A função afim que liga os pontos $(0, 0)$ e $(1,  1)$ é $f(x)=x$

  A função afim que liga os pontos $(0, 0)$ e $(2, -1)$ é $g(x)=-\frac{x}{2}$

  A função afim que liga os pontos $(1, 1)$ e $(2, -1)$ é $p(x)=3-2x$

  \textbf{Primeira parte $0\leq x\leq 1$ }
  \begin{align*}
    r(x) & = x + 1                                                     \\
    h(x) & = f(x) - g(x) = x -\left(-\frac{x}{2}\right) = \frac{3}{2}x \\
    a    & = 0                                                         \\
    b    & = 1
  \end{align*}
  \vspace{-2\baselineskip}
  \begin{align*}
    V_1 & = 2\pi \int_a^b r(x) h(x) \, dx                               \\
        & = 2\pi \int_0^1 (x+1) \frac{3x}{2} \, dx                      \\
        & = 3\pi\int_0^1 x^2+x \, dx                                    \\
        & = 3\pi\left(\frac{x^3}{3} + \frac{x^2}{2}\right)\evalat{0}{1} \\
        & = 3\pi \left( \frac{1}{3} + \frac{1}{2}\right)                \\
        & = \frac{5\pi}{2}
  \end{align*}

  \textbf{Segunda parte $1<x\leq 2$}
  \begin{align*}
    r(x) & = x + 1           \\
    h(x) & = p(x) - g(x)
           = \left(3-2x\right) -\left(-\frac{x}{2}\right)
           = 3 -2x +\frac{x}{2}
           = 3\left(1 -\frac{x}{2}\right) \\
    a    & = 1               \\
    b    & = 2
  \end{align*}
  \vspace{-2\baselineskip}
  \begin{align*}
    V_2 & = 2\pi \int_a^b r(x) h(x) \, dx                               \\
        & = 2\pi \int_1^2 (x+1) 3\left(1 -\frac{x}{2}\right) \, dx      \\
        & = 6\pi \int_1^2 x -\frac{x^2}{2} + 1 -\frac{x}{2}\, dx        \\
        & = 6\pi \int_1^2 1 + \frac{x}{2} -\frac{x^2}{2} \, dx          \\
        & = 6\pi \left(x +\frac{x^2}{4} - \frac{x^3}{6}\right)\evalat{1}{2}  \\
        & = 6\pi \left[
                 \left(2 +\frac{2^2}{4} - \frac{2^3}{6}\right)
               - \left(1 +\frac{1^2}{4} - \frac{1^3}{6}\right)
                 \right]  \\
        & = 6\pi \left(
                 2 + 1 - \frac{4}{3}
                -1 -\frac{1}{4} + \frac{1}{6}
                 \right)  \\
        & = 6\pi \,\frac{24 -16 -3 + 2}{12}  \\
        & = 6\pi \,\frac{7}{12} \\
        & = \frac{7\pi}{2}
  \end{align*}


  \textbf{Volume total}
  \[
    V
    = V_1 + V_2
    = \frac{5\pi}{2} + \frac{7\pi}{2}
    = 6\pi
  \]

\end{answer}

% 02
%------------------------------------------------------------------------------%
%\vspace{\baselineskip}
\exercise{20}
Calcule a integral
\(
  \int \frac{\sin(x)}{\,\cos^2(x)- 3\cos(x)\,}\, dx
\)

\begin{answer}
  \[
    F = \int \frac{\sin(x)}{\,\cos^2(x)- 3\cos(x)\,}\, dx
  \]
  \[
    u = \cos(x)
    \qquad\qquad
    du = -\sin(x)du
  \]
  \[
    F
    = \int \frac{-du}{u^2-3u}
    = \int \frac{du}{3u - u^2}
  \]
  Por frações parciais temos
  \[
    \frac{1}{3u-u^2}
    = \frac{1}{u(3-u)}
    = \frac{a}{u} + \frac{b}{3-u}
  \]
  \[
    1 = a(3-u) + bu = 3a + u(b-a)
  \]
  \[
    3a = 1 \qquad b-a = 0
  \]
  \[
    a = \frac{1}{3} \qquad b = a = \frac{1}{3}
  \]
  Portanto
  \begin{align*}
    F
    & = \int \frac{\sfrac{1}{3}}{u} + \frac{\sfrac{1}{3}}{3-u} du \\
    & = \frac{1}{3}\int \frac{du}{u} + \frac{1}{3}\int \frac{du}{3-u} \\
    & = \frac{1}{3}\ln\abs{u} - \frac{1}{3}\ln\abs{3-u} + c \\
    & = \frac{1}{3}\big(\ln\abs{\cos(x)} - \ln\abs{3-\cos(x)}\big) + c \\
    & = \ln\sqrt[3]{\abs{\frac{\cos(x)}{3-\cos(x)}}} + c
  \end{align*}

\end{answer}

% 03
%------------------------------------------------------------------------------%
%\vspace{\baselineskip}
\exercise{20}
Use o teste da integal para analisar a convergência da série
\(
  \sum_{n=1}^\infty \frac{e^n}{1+e^{2n}}
\)

\begin{answer}
  Vamos usar o teste da integral com a função \(f(x) = \frac{e^x}{1+e^{2x}}\)

  A função $f$ é sempre positiva

  A função $f$ é contínua para todo $x$ real

  Derivando $f$ temos
  \begin{align*}
    f'(x)
    & = \frac{d}{dx} \left(\frac{e^x}{1+e^{2x}}\right) \\
    & = \frac{\big(e^x\big)'\big(1+e^{2x}\big) - e^x\big(1+e^{2x}\big)'}{\big(1+e^{2x}\big)^2} \\
    & = \frac{e^x+e^xe^{2x} - e^xe^{2x}2}{\big(1+e^{2x}\big)^2} \\
    & = \frac{e^x +e^{3x} - 2e^{3x}}{\big(1+e^{2x}\big)^2} \\
    & = \frac{e^x - e^{3x}}{\big(1+e^{2x}\big)^2} \\
    & = \frac{e^x}{\big(1+e^{2x}\big)^2}\left(1 - e^{2x}\right)
  \end{align*}
  A derivada será negativa sempre que \(e^{2x}>1\), ou seja, sempre que \(x>0\)

  Calculando a primitiva
  \begin{align*}
    F(x)
    & = \int \frac{e^x}{1+e^{2x}} \, dx \qquad\qquad u = e^x \qquad du = e^x\,dx \\
    & = \int \frac{du}{1+u^2} \\
    & = \arctan\left(u\right) + c \\
    & = \arctan\left(e^x\right) + c
  \end{align*}
  Calculando a integral imprópria
  \begin{align*}
    \int_1^\infty \frac{e^x}{1+e^{2x}} \,dx
    & = \lim_{b\to\infty} \int_1^b \frac{e^x}{1+e^{2x}} dx \\
    & = \lim_{b\to\infty} \left(\arctan\left(e^x\right) \evalat{1}{b}\right) \\
    & = \lim_{b\to\infty} \left(\arctan\left(e^b\right) - \arctan\left(e\right)\right)
    \qquad\qquad p = e^b \\
    & = \lim_{p\to\infty} \left(\arctan\left(p\right) - \arctan\left(e\right)\right) \\
    & = \frac{\pi}{2} - \arctan\left(e\right)
  \end{align*}
  Como a integral converge a série também converge.
\end{answer}

% 04
%------------------------------------------------------------------------------%
%\vspace{\baselineskip}
\exercise{20}
Sabendo que \(f(1)=0\) e que, para \(n=1, 2, 3, \dots\),
\[
  \frac{d^n}{dx^n}f(x) = \frac{(-1)^n(n-1)!}{(3-x)^n}
\]

\begin{enumerate}[label=\alph*)]
  \item\ [04] Construa a Série de Taylor de $f(x)$ centrada em 1
  \item\ [04] Determine o intervalo aberto onde a série converge absolutamente
  \item\ [06] Use o polinômio de Taylor de terceiro grau para aproximar \(f(2)\)
  \item\ [06] Estime o erro cometido na aproximação
\end{enumerate}
Não é necessário encontrar a representação decimal dos valores calculados,
mas é necessário realizar as simplificações elementares.

\begin{answer}
\end{answer}

%------------------------------------------------------------------------------%
\end{document}
%------------------------------------------------------------------------------%
